%=======================================================================
\chapter{Construction des ensembles de caractéristiques (Phase~2)}
\label{ch:ensembles}

\section{Construction des ensembles de caractéristiques}

Après l'identification des paramètres optimaux lors de la Phase~1, la deuxième phase du
projet consiste à construire plusieurs ensembles de caractéristiques afin d'étudier l'apport
individuel et combiné des différents descripteurs.

L'objectif est d'évaluer dans quelle mesure les informations statistiques, texturales,
fréquentielles et structurelles sont complémentaires pour la classification des tumeurs
cérébrales. Pour cela, plusieurs configurations de fusion ont été définies puis comparées à
travers un benchmark systématique des modèles de Machine Learning.

Les ensembles étudiés sont présentés dans le Tableau~\ref{tab:feature-sets}.

\begin{table}[htbp]
  \centering
  \caption{Ensembles de caractéristiques évalués.}
  \label{tab:feature-sets}
  \begin{tabular}{@{}cl@{}}
    \toprule
    \thd{Nom} & \thd{Contenu} \\
    \midrule
    A     & Caractéristiques statistiques \\
    B     & Texture (GLCM + LBP optimisés) \\
    C     & DWT optimisé \\
    D     & HOG optimisé \\
    A+B   & Statistiques + Texture \\
    B+C   & Texture + DWT \\
    A+B+C & Statistiques + Texture + DWT \\
    Full  & Statistiques + Texture + DWT + HOG \\
    \bottomrule
  \end{tabular}
\end{table}

Chaque ensemble représente un niveau différent de complexité et de richesse descriptive. Les
ensembles simples permettent d'évaluer la contribution individuelle de chaque famille de
caractéristiques, tandis que les ensembles fusionnés permettent d'analyser leur
complémentarité.

La Figure~\ref{fig:phase2-pipeline} illustre le pipeline de construction des huit ensembles
à partir des branches A, B, C et D.

\begin{figure}[htbp]
  \centering
  \begin{tikzpicture}[
    node distance=0.55cm and 0.35cm,
    box/.style={rectangle, rounded corners=2pt, draw=accent, fill=headbg,
                minimum height=0.65cm, align=center, font=\small},
    fused/.style={rectangle, rounded corners=2pt, draw=accent2, fill=headbg,
                  minimum height=0.65cm, align=center, font=\small},
    arr/.style={-{Stealth[length=2mm]}, thick, accent}
  ]
    \node[box] (a) {A\\Stats (11)};
    \node[box, right=of a] (b) {B\\GLCM+LBP (20)};
    \node[box, right=of b] (c) {C\\DWT (16)};
    \node[box, right=of c] (d) {D\\HOG (1944)};
    \node[fused, below=1.2cm of b] (ab) {A+B};
    \node[fused, below=1.2cm of c] (bc) {B+C};
    \node[fused, below=2.4cm of b] (abc) {A+B+C};
    \node[fused, below=3.6cm of b] (full) {Full (1991)};
    \draw[arr] (a) -- (ab);
    \draw[arr] (b) -- (ab);
    \draw[arr] (b) -- (bc);
    \draw[arr] (c) -- (bc);
    \draw[arr] (ab) -- (abc);
    \draw[arr] (c) -- (abc);
    \draw[arr] (abc) -- (full);
    \draw[arr] (d) |- (full);
  \end{tikzpicture}
  \caption{Pipeline de construction des ensembles de caractéristiques (Phase~2).}
  \label{fig:phase2-pipeline}
\end{figure}

\subsection{Ensemble A : Caractéristiques statistiques}

Cet ensemble regroupe uniquement les \dimStats{} mesures statistiques globales extraites des
images IRM~: moyenne, écart-type, variance, entropie, asymétrie, kurtosis et les percentiles
d'intensité (p10 à p90, dont p50 = médiane).

Ces caractéristiques fournissent une représentation compacte de la distribution des niveaux
de gris mais ne prennent pas en compte l'organisation spatiale des pixels.

\subsection{Ensemble B : Texture (GLCM + LBP)}

L'ensemble B combine les informations texturales obtenues à partir des descripteurs GLCM
(\dimGLCM{} propriétés) et LBP (\dimLBP{} bins, $P=12$) optimisés en Phase~1, soit
\dimTexture{} dimensions au total.

Cette représentation vise à capturer les relations spatiales entre pixels ainsi que les
microstructures locales présentes dans les tissus cérébraux.

\subsection{Ensemble C : DWT}

Cet ensemble contient exclusivement les \dimDWT{} caractéristiques fréquentielles extraites
par Transformée en Ondelettes Discrètes (ondelette \texttt{db2}, niveau~1).

L'objectif est d'étudier la capacité des informations multi-résolutions à discriminer les
différentes catégories de tumeurs.

\subsection{Ensemble D : HOG}

L'ensemble D regroupe les \dimHOG{} caractéristiques issues du descripteur HOG optimisé.

Ces caractéristiques représentent principalement les contours, les gradients et les
structures géométriques présentes dans les images.

\subsection{Ensembles fusionnés}

Afin d'exploiter la complémentarité des différents descripteurs, plusieurs stratégies de
fusion ont été étudiées~:
\begin{itemize}[leftmargin=1.4em,itemsep=2pt]
  \item \textbf{A+B}~: combinaison des informations statistiques et texturales~;
  \item \textbf{B+C}~: combinaison des textures et des informations fréquentielles~;
  \item \textbf{A+B+C}~: fusion des statistiques, textures et ondelettes~;
  \item \textbf{Full}~: fusion de l'ensemble des caractéristiques disponibles
        (\dimFull{} dimensions brutes).
\end{itemize}

Cette approche permet d'évaluer l'impact de l'ajout progressif d'informations complémentaires
sur les performances de classification.

% -----------------------------------------------------------------------
\section{Nettoyage et préparation des caractéristiques}

La fusion de plusieurs familles de descripteurs conduit à la génération de vecteurs de grande
dimension. Une telle représentation peut contenir des variables redondantes ou peu
informatives susceptibles de dégrader les performances des modèles.

Afin d'améliorer la qualité des données et de réduire le risque de surapprentissage, une
procédure de nettoyage a été appliquée à chaque ensemble de caractéristiques avant
l'entraînement des modèles.

\subsection{Suppression des caractéristiques à faible variance}

La première étape consiste à appliquer un filtre basé sur la variance
(\emph{Variance Threshold}, seuil~$=0$).

Les variables dont la variance est proche de zéro apportent généralement peu d'information
discriminante puisqu'elles prennent pratiquement la même valeur pour l'ensemble des
observations.

La suppression de ces caractéristiques permet~:
\begin{itemize}[leftmargin=1.4em,itemsep=2pt]
  \item de réduire la dimension des données~;
  \item d'éliminer les variables non informatives~;
  \item d'accélérer les étapes ultérieures de traitement.
\end{itemize}

Cette opération constitue un premier niveau de sélection automatique des caractéristiques.

\subsection{Élimination des caractéristiques fortement corrélées}

La présence de variables fortement corrélées introduit une redondance importante dans la
représentation des données.

Pour limiter ce phénomène, une matrice de corrélation est calculée sur l'ensemble des
caractéristiques. Lorsqu'une paire de variables présente une corrélation absolue supérieure
au seuil $\rho = \corrThreshold{}$, l'une des deux variables est supprimée.

Cette étape permet~:
\begin{itemize}[leftmargin=1.4em,itemsep=2pt]
  \item de réduire la multicolinéarité~;
  \item d'améliorer la stabilité des modèles~;
  \item de diminuer la complexité des ensembles fusionnés.
\end{itemize}

La suppression des variables redondantes est particulièrement importante pour les ensembles
comportant plusieurs milliers de caractéristiques, comme HOG ou Full.

\subsection{Standardisation des caractéristiques}

Les différentes familles de descripteurs produisent des variables dont les échelles peuvent
être très différentes. Par exemple, certaines caractéristiques statistiques prennent des
valeurs comprises entre 0 et 1, tandis que certaines caractéristiques HOG ou DWT peuvent
atteindre des amplitudes beaucoup plus élevées.

Afin d'assurer une contribution équilibrée de chaque variable lors de l'apprentissage, une
standardisation (\texttt{StandardScaler}) est appliquée à l'ensemble des caractéristiques
après nettoyage. Chaque variable est transformée selon~:
\begin{equation}
  z = \frac{x-\mu}{\sigma}
\end{equation}
où $x$ représente la valeur initiale, $\mu$ la moyenne et $\sigma$ l'écart-type estimés sur
l'ensemble d'apprentissage.

Cette transformation produit des variables centrées et réduites, facilitant l'apprentissage
des modèles sensibles à l'échelle des données tels que SVM, KNN, Logistic Regression et MLP.

% -----------------------------------------------------------------------
\section{Réduction de la redondance et analyse des dimensions}

L'une des conséquences directes de la fusion de plusieurs descripteurs est l'augmentation
significative du nombre de variables. Une analyse de la réduction dimensionnelle obtenue
après le nettoyage a donc été réalisée.

Pour chaque ensemble de caractéristiques, les éléments suivants ont été étudiés~:
\begin{itemize}[leftmargin=1.4em,itemsep=2pt]
  \item nombre initial de caractéristiques~;
  \item nombre de caractéristiques supprimées~;
  \item nombre final de caractéristiques conservées~;
  \item ratio de conservation.
\end{itemize}

Le ratio de conservation est défini par~:
\begin{equation}
  \mathrm{Ratio}_{\mathrm{conservation}}
  = \frac{N_{\mathrm{conservées}}}{N_{\mathrm{initial}}}
  \times 100
\end{equation}

Cette mesure permet d'évaluer l'impact du processus de nettoyage sur chaque ensemble.

Le Tableau~\ref{tab:cleaning-dims} et les Figures~\ref{fig:cleaning-kept} et
\ref{fig:cleaning-dims} synthétisent les résultats obtenus.

\begin{table}[htbp]
  \centering
  \caption{Dimensions avant et après nettoyage pour chaque ensemble.}
  \label{tab:cleaning-dims}
  \small
  \begin{tabular}{@{}lrrrrr@{}}
    \toprule
    \thd{Ensemble} & \thd{Brut} & \thd{Après var.} & \thd{Nettoyé} & \thd{Supprimé} &
    \thd{Conservé (\%)} \\
    \midrule
    A     & 11          & 11          & 9           & 2    & 81,8 \\
    B     & 20          & 20          & 18          & 2    & 90,0 \\
    C     & 16          & 16          & 12          & 4    & 75,0 \\
    D     & 1944        & 1944        & 490         & 1454 & 25,2 \\
    A+B   & 31          & 31          & 27          & 4    & 87,1 \\
    B+C   & 36          & 36          & 30          & 6    & 83,3 \\
    A+B+C & 47          & 47          & 36          & 11   & 76,6 \\
    Full  & 1991        & 1991        & 526         & 1465 & 26,4 \\
    \bottomrule
  \end{tabular}
\end{table}

Les résultats montrent que les ensembles de grande dimension, notamment ceux incluant HOG et
les combinaisons fusionnées, présentent une proportion importante de caractéristiques
redondantes. À l'inverse, l'ensemble statistique A possède naturellement une dimension
réduite et nécessite peu de filtrage.

\begin{figure}[htbp]
  \centering
  \includegraphics[width=0.92\textwidth]{feature_cleaning_kept_ratio}
  \caption{Ratio de conservation des caractéristiques après nettoyage.}
  \label{fig:cleaning-kept}
\end{figure}

\begin{figure}[htbp]
  \centering
  \includegraphics[width=\textwidth]{feature_cleaning_dims}
  \caption{Comparaison des dimensions des ensembles avant et après réduction.}
  \label{fig:cleaning-dims}
\end{figure}

L'application conjointe du \emph{Variance Threshold} et de l'élimination des corrélations
permet ainsi de conserver uniquement les caractéristiques les plus informatives tout en
réduisant significativement la complexité des données. L'ensemble Full passe par exemple de
1991 à \cleanFull{} variables (soit 26,4\% conservées).

Cette étape constitue un élément essentiel du pipeline, car elle améliore à la fois la
robustesse des modèles, leur temps d'entraînement et leur capacité de généralisation.

\begin{figure}[htbp]
  \centering
  \includegraphics[width=0.88\textwidth]{tsne_best_feature_set}
  \caption{Projection t-SNE (2D) de l'ensemble Full(opt) nettoyé (\cleanFull{}
    caractéristiques) sur l'ensemble d'apprentissage.}
  \label{fig:tsne-full}
\end{figure}

La Figure~\ref{fig:tsne-full} illustre la séparabilité visuelle des trois classes dans
l'espace des caractéristiques de l'ensemble fusionné complet après nettoyage.

\subsection{Synthèse de la Phase~2}

À l'issue de cette phase, huit ensembles de caractéristiques nettoyés et standardisés sont
obtenus. Ces représentations serviront de base au benchmark des algorithmes de Machine
Learning présenté au chapitre suivant.

L'objectif sera alors d'évaluer l'impact de chaque famille de descripteurs ainsi que
l'intérêt des stratégies de fusion sur les performances de classification des tumeurs
cérébrales.
